\bookchapter{Annexe H : feuille de route des versions}

\chapterintro{Ce livre continue d'être mis à jour. Voici les orientations envisagées pour les versions à venir ; dis-nous aussi dans tes retours ce que tu aimerais le plus y voir.}

\textbf{v1.8 est en place : ajout des septième, huitième et neuvième parties « Pratique de la programmation par IA », « Pratique du développement Android » et « Pratique du développement iOS » (chapitres 37-69, soit 33 chapitres)}\par

De la révolution de paradigme de la programmation par IA, les cinq principes de pensée, l'ingénierie du contexte et des besoins, le panorama des outils et modèles de 2026, la pratique des API, l'ingénierie de la conversation, trois projets complets, l'assurance qualité et la montée en compétence — pour te faire passer de « savoir utiliser l'IA » à « diriger l'IA qui code » ; jusqu'à la chaîne complète de développement des deux plateformes Android et iOS : installation de l'environnement (Android Studio / Xcode), langage en accéléré (Kotlin / Swift), frameworks d'interface (Jetpack Compose / SwiftUI), données et réseau, projets pratiques (carnet de comptes + météo), tests et publication (Google Play / App Store) et aide-mémoire — pour te faire grandir d'utilisateur numérique jusqu'à développeur d'app.

\begin{tcbinfo}{Orientations prévues pour la prochaine version (v1.9)}
\begin{itemize}
  \item Téléphone et objets portables : guide de synergie approfondie avec l'Apple Watch, les montres Huawei et les bagues connectées (Oura/Galaxy Ring)
  \item Voyage et téléphone : eSIM en pratique, solutions de data à l'étranger, cartes hors ligne, paiements à l'étranger, comparatif des apps de traduction
  \item Téléphone et santé : sport et bien-être, suivi du sommeil, rythme cardiaque et SpO2, santé féminine, apps de santé mentale
  \item Vidéos courtes créées par IA : tutoriels pratiques de CapCut IA, Kling IA, humains numériques et doublage IA
  \item Initiation aux mini-séries verticales : structure de scénario, techniques de tournage, règles des plateformes
\end{itemize}

\end{tcbinfo}

\begin{tcbwarn}{Remarque sur l'actualité des informations}
Les technologies du téléphone évoluent extrêmement vite ; ce livre est rédigé d'après la situation technique et du marché de la mi-2026. Les capacités IA, les modèles de puces, les versions des systèmes et les fourchettes de prix ont peut-être déjà changé au moment de ta lecture. Pour les étapes concrètes, fie-toi à la version système réelle de ton appareil ; pour tes décisions d'achat, consulte les derniers tests et les retours d'utilisateurs.

\end{tcbwarn}
